\documentclass[11pt,a4paper]{article}
\usepackage[margin=2.5cm]{geometry}
\usepackage[T1]{fontenc}
\usepackage[utf8]{inputenc}
\usepackage{amsmath,amssymb}
\usepackage{amsthm}     % theorem environment used in the chapter example
\newtheorem{theorem}{Theorem}
\usepackage{multicol}   % columns, used in the poster-style example
\usepackage{enumitem}   % control list spacing
\usepackage{booktabs}   % table rules in the transliteration table

% ---------------------------------------------------------------
% MODULE V - Open Ended
% Syllabus unit list: 1 Installation, 2 Overleaf, 3 Writing a chapter,
% 4 Beamer slides and posters, 5 Transliteration (Sanskrit, Hindi,
% Malayalam). The main compiled example uses pdflatex. Items that need
% XeLaTeX or LuaLaTeX are marked and shown as listings only.
% ---------------------------------------------------------------

\title{Module V: Open-Ended Projects in \LaTeX}
\author{Study Notes -- MAT5FS112}
\date{}

\begin{document}
\maketitle
\tableofcontents

% ===============================================================
\section{Installing LaTeX}
% ===============================================================

You need a \emph{distribution} (the engine plus packages) and an
editor. The usual choices are:
\begin{itemize}
  \item \textbf{TeX Live} (Linux, macOS, Windows). On Arch Linux, install
    \texttt{texlive-basic texlive-latex texlive-latexrecommended
    texlive-latexextra texlive-fontsrecommended texlive-bibtexextra}
    with \texttt{pacman}. On Debian or Ubuntu: \texttt{texlive-full}
    (large, but complete).
  \item \textbf{MiKTeX} (Windows), which installs missing packages on
    demand.
  \item \textbf{MacTeX} (macOS).
\end{itemize}

After installation, check the tools are on the path:
\begin{verbatim}
pdflatex --version
bibtex --version
makeindex --version
\end{verbatim}

A good editor (TeXstudio, TeXworks, VS Code with LaTeX Workshop, or
Vim with a LaTeX plugin) builds the document and shows errors.

% ===============================================================
\section{Overleaf}
% ===============================================================

Overleaf (\texttt{overleaf.com}) is an online editor. It needs no
installation, has collaborative editing and compiles the document on
its servers. The steps are:
\begin{enumerate}[noitemsep]
  \item Create a free account and choose \emph{New Project}, then
    \emph{Blank Project}.
  \item Paste the code from any file in these notes into \texttt{main.tex}.
  \item Press \emph{Recompile}. The PDF appears on the right, and the
    compile log is under \emph{Logs and output}.
  \item Use the menu \emph{Menu $\rightarrow$ Settings} to choose the
    compiler (pdfLaTeX, XeLaTeX or LuaLaTeX) and the main file.
\end{enumerate}

Keep a local copy of your \texttt{.tex} file. Overleaf can be unavailable,
and the same code compiles on a local TeX Live.

% ===============================================================
\section{Writing a chapter of a book}
% ===============================================================

A book uses the \texttt{report} class (or \texttt{book}), where the top
level is \verb|\chapter|. The skeleton below is a template for a chapter
of a textbook with theorems and equations:

\begin{verbatim}
\documentclass[11pt,a4paper]{report}
\usepackage{amsmath,amsthm,amssymb}
\newtheorem{theorem}{Theorem}[chapter]
\begin{document}
\chapter{Introduction to Sets}
\section{Basic notions}
\begin{theorem}
...
\end{theorem}
\end{document}
\end{verbatim}

Here is the same structure, compiled, in this \texttt{article} file.
\begin{theorem}[Sample result]
Every finite set has finitely many subsets, namely $2^n$ for a set with
$n$ elements.
\end{theorem}

\begin{proof}
Each element is either in a subset or not, which gives two choices per
element and $2^n$ choices in all.
\end{proof}

% ===============================================================
\section{Slides with Beamer}
% ===============================================================

\texttt{beamer} is a document class for presentations. Each slide is a
\verb|\frame|. The code below belongs in a file of its own that starts
with \verb|\documentclass{beamer}|; it is shown here as a listing:

\begin{verbatim}
\documentclass{beamer}
\usetheme{Madrid}
\usecolortheme{beaver}
\title{Linear Algebra Basics}
\author{Your Name}
\date{\today}
\begin{document}
\frame{\titlepage}
\begin{frame}{Vector spaces}
  \begin{block}{Definition}
    A vector space is a set closed under addition and scalar multiplication.
  \end{block}
  \pause
  \begin{itemize}
    \item Eight axioms must hold.
    \item Examples: $\mathbb{R}^n$, polynomials.
  \end{itemize}
\end{frame}
\end{document}
\end{verbatim}

Key ideas:
\begin{itemize}
  \item \verb|\frame{...}| or \verb|\begin{frame}{Title}...\end{frame}|.
  \item \verb|\pause| reveals the content in steps.
  \item \verb|\begin{block}{...}| gives a highlighted box.
\end{itemize}

% ===============================================================
\section{Posters}
% ===============================================================

A poster is a single large page with several columns. The
\texttt{beamerposter} package is the usual tool. Without it, a
poster-like layout can be built with \texttt{multicols} and a
\texttt{minipage}. Here is a compiled version of a two-column poster
section:

\begin{multicols}{2}
\section*{Motivation}
Matrices describe linear maps. Their determinant measures how the
map scales area:
\[
  \det\begin{pmatrix} a & b \\ c & d \end{pmatrix} = ad - bc.
\]

\section*{Result}
A matrix is invertible if and only if its determinant is non-zero.
\columnbreak
\section*{Example}
For $A = \begin{pmatrix} 2 & 1 \\ 1 & 3 \end{pmatrix}$, $\det A = 5 \neq 0$,
so $A$ is invertible.
\end{multicols}

% ===============================================================
\section{Transliteration of Indian languages}
% ===============================================================

Transliteration writes an Indian script in Latin letters. Scholarly
Sanskrit uses IAST (International Alphabet of Sanskrit Transliteration),
whose letters need diacritics. In \texttt{pdflatex} these are made
with accent commands. This is compiled in this file:

\begin{center}
\begin{tabular}{ll}
\toprule
Command & Letter \\
\midrule
\verb|\={a}| & $\bar{a}$ as in \emph{R\={a}ma} \\
\verb|\'{s}| & \'{s} as in \emph{\'{s}iva} \\
\verb|\d{t}| & \d{t} as in \emph{\d{t}ha} \\
\verb|\d{n}| & \d{n} as in \emph{Ka\d{n}a} \\
\verb|\.{m}| & \.{m} as in \emph{Sa\.{m}sk\d{r}ta} \\
\verb|\d{r}| & \d{r} as in \emph{\d{r}\d{s}i} \\
\bottomrule
\end{tabular}
\end{center}

Example in running text: \emph{R\={a}m\={a}yana} and
\emph{Bhagavad-g\={\i}t\={a}}.

\subsection{Devanagari and Malayalam (XeLaTeX or LuaLaTeX)}

Printing the script itself needs a Unicode engine and a font that has
the glyphs. The usual setup is \texttt{fontspec} with
\texttt{polyglossia}. This listing is \emph{not} compiled here, because
it needs XeLaTeX and an installed font. Replace each placeholder with
the word typed in a Unicode-capable editor:

\begin{verbatim}
\documentclass{article}
\usepackage{fontspec}
\usepackage{polyglossia}
\setdefaultlanguage{hindi}              % or sanskrit, malayalam
\newfontfamily\devanagarifont{Noto Serif Devanagari}
\newfontfamily\malayalamfont{Noto Serif Malayalam}
\begin{document}
{\devanagarifont <Devanagari word>} \quad {\malayalamfont <Malayalam word>}
\end{document}
\end{verbatim}

Compile with \verb|xelatex file.tex| (or \verb|lualatex|). Check the font
names with \verb|fc-list| first, because they must match the installed
font exactly.

% ===============================================================
\section{Checklist for an open-ended project}
% ===============================================================

\begin{enumerate}
  \item Set up the distribution and the editor, and compile a test file.
  \item Choose the document class (\texttt{article}, \texttt{report} or
    \texttt{beamer}).
  \item Add packages only as needed, and keep the preamble short.
  \item Compile twice for references, the table of contents and the
    list of figures.
  \item Run BibTeX or makeindex when the document uses them.
  \item Read the log: the first error line is the one to fix.
\end{enumerate}

\end{document}
